% Detailed solutions for the 24 canonical VI/41 exercises.

% VI41-EXERCISE-SOLUTION-01
\paragraph{Solution to Exercise VI.41.01.}
We verify the three axioms. Reflexivity follows from \(D-D=\operatorname{div}(1)\). If
\(D-E=\operatorname{div}(f)\), then \(E-D=\operatorname{div}(f^{-1})\), proving symmetry.
If \(D-E=\operatorname{div}(f)\) and \(E-F=\operatorname{div}(g)\), then
\[
D-F=\operatorname{div}(fg),
\]
so transitivity holds. Thus linear equivalence is an equivalence relation.

% VI41-EXERCISE-SOLUTION-02
\paragraph{Solution to Exercise VI.41.02.}
Since \(\Cl(X)=\Div(X)/\Prin(X)\), the zero class is the coset \(\Prin(X)\).
Therefore \([D]=0\) iff \(D\in\Prin(X)\), i.e. iff \(D=\operatorname{div}(f)\) for some
\(f\in K(X)^\times\).

% VI41-EXERCISE-SOLUTION-03
\paragraph{Solution to Exercise VI.41.03.}
The ring \(k[t]\) is a PID, hence a normal Noetherian UFD. Every height-one prime is
generated by an irreducible polynomial \(p(t)\), and its prime divisor equals
\(\operatorname{div}(p)\). Hence every Weil divisor is principal and
\[
\Cl(\Spec k[t])=0.
\]

% VI41-EXERCISE-SOLUTION-04
\paragraph{Solution to Exercise VI.41.04.}
Affine \(n\)-space has coordinate ring \(k[x_1,\ldots,x_n]\), a polynomial UFD.
The UFD criterion for normal Noetherian domains therefore gives
\[
\Cl(\AA^n_k)=0.
\]
Equivalently, every prime divisor is an irreducible hypersurface and is principal.

% VI41-EXERCISE-SOLUTION-05
\paragraph{Solution to Exercise VI.41.05.}
Let \(\mathfrak p\) have height one. Since \(\Cl(A)=0\), choose
\(f\in\Frac(A)^\times\) with \(\operatorname{div}(f)=[\mathfrak p]\).
All height-one valuations of \(f\) are nonnegative, so the Krull-domain identity
\[
A=\bigcap_{\height\mathfrak q=1}A_{\mathfrak q}
\]
implies \(f\in A\). If \(a\in\mathfrak p\), then
\(v_{\mathfrak p}(a/f)\ge0\), and for every other height-one prime
\(v_{\mathfrak q}(a/f)=v_{\mathfrak q}(a)\ge0\). Thus \(a/f\in A\), so
\(\mathfrak p\subseteq(f)\). Since \(v_{\mathfrak p}(f)=1\), also \((f)\subseteq\mathfrak p\).
Hence \(\mathfrak p=(f)\).

% VI41-EXERCISE-SOLUTION-06
\paragraph{Solution to Exercise VI.41.06.}
If \(A\) is a UFD, then \(\Cl(A)=0\). The localization exact sequence gives a
surjection
\[
\Cl(A)\twoheadrightarrow\Cl(S^{-1}A),
\]
so \(\Cl(S^{-1}A)=0\). Since localization preserves normality and Noetherianity,
the UFD criterion implies that \(S^{-1}A\) is again a UFD.

% VI41-EXERCISE-SOLUTION-07
\paragraph{Solution to Exercise VI.41.07.}
For \(A=k[x,y,z]/(xy-z^2)\),
\[
A/(x,z)\cong k[y],\qquad A/(y,z)\cong k[x].
\]
Both quotients have dimension \(1\), whereas \(\dim A=2\). Thus
\[
\height(x,z)=\height(y,z)=1,
\]
so both ideals define prime Weil divisors.

% VI41-EXERCISE-SOLUTION-08
\paragraph{Solution to Exercise VI.41.08.}
At \(L=V(x,z)\), the element \(y\) is a unit and
\[
x=z^2/y.
\]
The height-one localization is a DVR with \(z\) as a uniformizer, so
\(v_L(z)=1\) and \(v_L(x)=2\). Since \(\sqrt{(x)}=(x,z)\), there is no other
height-one component of \(V(x)\). Therefore
\[
\operatorname{div}(x)=2L.
\]

% VI41-EXERCISE-SOLUTION-09
\paragraph{Solution to Exercise VI.41.09.}
Setting \(z=0\) in \(xy=z^2\) gives \(xy=0\), so the codimension-one zero locus is
\(L\cup M\), where \(L=V(x,z)\) and \(M=V(y,z)\). At both generic points \(z\) is a
uniformizer, hence has order one. Thus
\[
\operatorname{div}(z)=L+M.
\]

% VI41-EXERCISE-SOLUTION-10
\paragraph{Solution to Exercise VI.41.10.}
Using \(\operatorname{div}(x)=2L\) and \(\operatorname{div}(z)=L+M\),
\[
\operatorname{div}(x/z)=2L-(L+M)=L-M.
\]
Hence \(L-M\) is principal, so \(L\sim M\) and \([L]=[M]\).

% VI41-EXERCISE-SOLUTION-11
\paragraph{Solution to Exercise VI.41.11.}
After inverting \(x\), solve the relation as \(y=z^2/x\). This gives inverse maps
between \(A_x\) and \(k[x,x^{-1},z]\), hence
\[
A_x\cong k[x,x^{-1},z].
\]
The right side is a localization of \(k[x,z]\), so it is a UFD.

% VI41-EXERCISE-SOLUTION-12
\paragraph{Solution to Exercise VI.41.12.}
Because \(A_x\) is a UFD, \(\Cl(A_x)=0\). A height-one prime disappears after
inverting \(x\) precisely when it contains \(x\). Since
\[
\sqrt{(x)}=(x,z),
\]
the unique such height-one prime is \(L\). Hence localization yields
\[
0=\Cl(A_x)\cong\Cl(A)/\langle[L]\rangle,
\]
so \([L]\) generates \(\Cl(A)\).

% VI41-EXERCISE-SOLUTION-13
\paragraph{Solution to Exercise VI.41.13.}
The cone ring is connected positively graded with \(A_0=k\). If \(uv=1\), compare
the highest nonzero homogeneous pieces of \(u\) and \(v\). Since \(A\) is a domain,
their product cannot cancel, so both highest degrees must be \(0\). Therefore
\(u,v\in k\), and
\[
A^\times=k^\times.
\]

% VI41-EXERCISE-SOLUTION-14
\paragraph{Solution to Exercise VI.41.14.}
In the UFD \(k[x,z]\), the irreducible \(x\) defines the \(x\)-adic valuation on
\(k(x,z)\). For \(c\in k^\times\),
\[
v_x(cx)=1.
\]
If \(cx=h^2\), then \(1=v_x(h^2)=2v_x(h)\), impossible. Thus \(cx\) is not a square.

% VI41-EXERCISE-SOLUTION-15
\paragraph{Solution to Exercise VI.41.15.}
We already know \([L]\) generates and \(2[L]=0\) from
\(\operatorname{div}(x)=2L\). If \([L]=0\), write \(L=\operatorname{div}(f)\). Then
\[
\operatorname{div}(f^2/x)=0.
\]
The Krull intersection formula implies \(f^2/x\in A^\times=k^\times\), so
\(f^2=cx\). Since \(K(A)=k(x,z)\), the previous exercise rules this out.
Hence \([L]\neq0\), so
\[
\Cl(A)\cong\mathbb Z/2\mathbb Z.
\]

% VI41-EXERCISE-SOLUTION-16
\paragraph{Solution to Exercise VI.41.16.}
The removed vertex has codimension two. Thus prime divisors of \(X\) and
\(U=X\setminus\{o\}\) correspond by intersection and closure. The function fields and
height-one local rings are unchanged, so principal divisors correspond as well. Hence
\[
\Cl(U)\cong\Cl(X)\cong\mathbb Z/2\mathbb Z.
\]

% VI41-EXERCISE-SOLUTION-17
\paragraph{Solution to Exercise VI.41.17.}
The punctured cone is smooth, so all its local rings are regular and hence UFDs:
\(U\) is locally factorial. But
\[
\Cl(U)\cong\mathbb Z/2\mathbb Z.
\]
Local factoriality only makes every Weil divisor locally principal (Cartier); it does
not force a single global rational function to generate each divisor.

% VI41-EXERCISE-SOLUTION-18
\paragraph{Solution to Exercise VI.41.18.}
Let \(U=D_+(x_0)\cong\AA^n\) and \(H=V_+(x_0)\). Since \(\Cl(U)=0\), the open-subset
exact sequence gives
\[
\mathbb Z[H]\twoheadrightarrow\Cl(\PP^n).
\]
Thus the hyperplane class generates the projective-space class group.

% VI41-EXERCISE-SOLUTION-19
\paragraph{Solution to Exercise VI.41.19.}
Write a rational function on \(\PP^n\) as \(F/G\), with \(F,G\) homogeneous of the
same degree after cancelling common factors. Its divisor is the zero divisor of \(F\)
minus that of \(G\). Their total degrees are equal, so
\[
\deg\operatorname{div}(F/G)=0.
\]

% VI41-EXERCISE-SOLUTION-20
\paragraph{Solution to Exercise VI.41.20.}
The hyperplane class generates. If \(m[H]=0\), then \(mH\) is principal, hence has
degree zero. But \(\deg(mH)=m\), so \(m=0\). Therefore the generator has infinite
order and
\[
\Cl(\PP^n_k)\cong\mathbb Z.
\]

% VI41-EXERCISE-SOLUTION-21
\paragraph{Solution to Exercise VI.41.21.}
Choose a linear form \(\ell\) not sharing a component with \(F\). Then
\[
\operatorname{div}(F/\ell^d)=V_+(F)-dV_+(\ell).
\]
Therefore
\[
[V_+(F)]=d[V_+(\ell)]=d[H],
\]
with all factor multiplicities included in \(V_+(F)\).

% VI41-EXERCISE-SOLUTION-22
\paragraph{Solution to Exercise VI.41.22.}
If \(X\setminus U\) has codimension at least two, no codimension-one irreducible
closed subset of \(X\) lies in the complement. Intersection with \(U\) therefore sends
prime divisors of \(X\) to prime divisors of \(U\). Conversely, closure in \(X\) sends
each prime divisor of \(U\) to a codimension-one irreducible closed subset of \(X\).
These operations are inverse, so
\[
\Div(X)\xrightarrow{\sim}\Div(U).
\]

% VI41-EXERCISE-SOLUTION-23
\paragraph{Solution to Exercise VI.41.23.}
The nonempty open \(U\) has the same function field as \(X\). For corresponding
prime divisors, the generic point and its local ring are the same in \(X\) and \(U\).
Thus
\[
\operatorname{ord}^{X}_Z(f)=\operatorname{ord}^{U}_{Z\cap U}(f)
\]
for every rational function \(f\). Hence principal divisors are preserved under the
divisor-group isomorphism.

% VI41-EXERCISE-SOLUTION-24
\paragraph{Solution to Exercise VI.41.24.}
On the cone,
\[
2L=\operatorname{div}(x),
\]
so \(2L\) is principal, hence Cartier. VI/40 proves that \(L\) itself is not Cartier at
the vertex because \((x,z)\) is not locally principal there. Therefore the least
positive Cartier multiple of \(L\) is \(2\), and \(L\) is \(\mathbb Q\)-Cartier of
index two.
